\chapter{Introduction}
\label{ch:introduction}

\section{Overview}

A great deal of useful knowledge lives inside documents. Companies keep their
policies, contracts and reports as files; researchers keep papers and notes;
students keep lecture material and textbooks; engineers keep specifications and
source code. These collections grow steadily, and after a while it becomes hard
to find a precise answer inside them, to see how the ideas in them connect, or
to turn a buried table into a chart that can be understood at a glance. The
information is present, but it is locked away in a form that rewards slow,
manual reading.

This thesis presents \emph{Synapse}, a web platform that helps a person or a
team make sense of their own collection of documents. A user connects a library
of files, and Synapse reads each file, understands its structure, and then lets
the user question it, visualise it and map it through a single application. The
answers it gives are grounded in the user's own material and carry a citation to
the exact passage they came from, so the user can check every claim. When the
library cannot support an answer, Synapse says so rather than inventing one.
Figure~\ref{fig:capabilities} summarises the four things the platform lets a
user do with a connected library.

\begin{figure}[H]
  \centering
  \includegraphics[width=0.92\textwidth]{fig_capabilities.png}
  \caption{The four capabilities Synapse provides over a connected document
  library. A user can ingest and understand a collection, ask grounded
  questions through a citing chat, create artefacts such as charts and
  documents through an agent, and map the corpus as an interactive knowledge
  graph. Teams share these libraries and collaborate in real time. All four
  capabilities act on the same processed library, so work done in one is
  available to the others.}
  \label{fig:capabilities}
\end{figure}

\section{Background and Motivation}

Large language models have made it natural to ask questions in plain language
and receive a fluent reply. On their own, however, these models suffer from two
well-known weaknesses. They only know what was present in their training data,
so they cannot speak about a private contract or an internal report that they
have never seen. And when they are unsure, they tend to answer anyway, producing
confident text that is not supported by any source. For everyday browsing this
is a minor nuisance. For someone deciding on the basis of a legal clause, a
financial figure or a medical guideline, an unsupported answer is worse than no
answer at all.

Retrieval-augmented generation was introduced to address the first weakness
\cite{lewis2020rag}. Instead of relying only on the model's memory, the system
first retrieves passages that are relevant to the question from a trusted
collection, and then asks the model to answer using those passages. This grounds
the reply in real material and makes it possible to show where each statement
came from. Retrieval-augmented generation has become the standard way to build a
question-answering system over a private corpus, and a large body of recent work
has refined how passages are retrieved, judged and combined
\cite{yan2024crag,asai2024selfrag,nguyen2025marag}.

Two practical problems remain, and they motivate this project. The first is that
real document collections are not made of clean paragraphs of text. They contain
scanned pages, spreadsheets, slide decks, figures, tables and source files. A
system that can only read born-digital prose will silently miss much of what a
collection actually holds. The second is that retrieving the right passage is
necessary but not sufficient. A useful system must also judge whether the
retrieved material truly supports an answer, decline to answer when it does not,
and present the result in whatever form the user needs, whether that is a short
written reply, a chart, a document or a picture of how the ideas connect. Synapse
is built around these two concerns.

\section{Problem Statement}

People and teams hold valuable knowledge in large, mixed collections of
documents, but they have no convenient and trustworthy way to use it. General
chatbots cannot read a private library, cannot ground their answers in it, and
will invent facts when the evidence is thin. Traditional search returns a list
of files rather than a direct, supported answer, and neither approach can read a
scanned page, draw a chart from a table, or reveal how the concepts in a
collection relate to one another. What is needed is a single system that ingests
a mixed document collection reliably, answers questions about it with honest,
cited, grounded replies, turns it into useful artefacts on request, and presents
its overall structure in a way a person can explore.

\section{Aim and Objectives}

The aim of this project is to design and build a practical, multi-tenant platform
that lets a user understand, question, visualise and map a private collection of
documents through one web application. To meet this aim, the following objectives
were set.

\begin{enumerate}[leftmargin=*,itemsep=2pt]
  \item To build an ingestion pipeline that reliably processes a mixed
        collection of documents, including born-digital files, scanned pages,
        spreadsheets and source code, into a searchable form.
  \item To design a retrieval-augmented chat that answers questions using only
        the user's library, cites the exact source of every claim, and abstains
        when the evidence is insufficient.
  \item To provide an agent that plans its work, asks for clarification when a
        request is genuinely ambiguous, and produces artefacts such as charts,
        written documents and images from the connected libraries.
  \item To extract the entities and relationships in a processed library and
        present them as an interactive knowledge graph the user can explore.
  \item To support teams, so that several users can share libraries and
        collaborate in real time.
  \item To evaluate the platform on an established document benchmark and on
        purpose-built datasets, and to report retrieval accuracy, answer
        quality, honesty and efficiency.
\end{enumerate}

\section{Scope of the Project}

The project covers the complete platform: the ingestion pipeline, the grounded
chat, the agent, the knowledge graph, team collaboration, and the supporting web
application and infrastructure. The system is multi-tenant, meaning that several
organisations can use it at once with their data kept separate. Documents are
processed automatically, and the user interacts only through the web
application.

The project does not aim to train new language or vision models; it uses
established hosted and open models and concentrates on the design that turns them
into a dependable product. It also does not attempt to be a general web agent
that operates arbitrary external websites; the agent works on the user's own
connected libraries and uploaded files. These boundaries kept the work focused on
the central goal of trustworthy document intelligence over a private collection.

\section{Project Contributions}

The main contributions of this work are as follows. First, it delivers a
single platform that brings together ingestion, grounded question answering,
artefact generation and knowledge mapping over a private library, where most
existing tools address only one of these. Second, it implements a corrective,
multi-agent retrieval loop that grades its own evidence and abstains rather than
guessing, and that attaches a citation to every claim. Third, it implements an
ingestion pipeline that handles scanned and structured documents, not only clean
prose, by combining layout detection, optical character recognition and a
vision-language model. Fourth, it provides an honest evaluation of the platform
on an established benchmark and on purpose-built datasets, reporting not only
accuracy but also how often the system is wrongly confident.

\section{Development Methodology}

The platform was built using an iterative and incremental software lifecycle, as
shown in Figure~\ref{fig:methodology}. Rather than attempting to specify and
build the whole system at once, the work was divided into a sequence of
increments, each adding one coherent piece of functionality such as the
ingestion pipeline, the chat, the agent or the knowledge graph. Every increment
passed through the same short cycle of planning, analysis, design,
implementation, testing and deployment, and produced a working part of the
system that could be tried and reviewed before the next increment began.

\begin{figure}[H]
  \centering
  \includegraphics[width=0.78\textwidth]{fig_methodology.png}
  \caption{The iterative and incremental lifecycle followed during development.
  Each feature of the platform passed through the same cycle of planning,
  analysis, design, implementation, testing and deployment. Every iteration
  delivered a working, tested increment, which kept the system usable
  throughout and allowed the design to be corrected early as the requirements
  became clearer.}
  \label{fig:methodology}
\end{figure}

This approach suited the project well. The most uncertain parts, such as how a
scanned page should be read or how the retrieval loop should decide to abstain,
benefited from being tried early on real data and refined in the light of what
worked. Because each increment was integrated into the running system, problems
in the way components fitted together were found while they were still cheap to
fix, and the platform remained demonstrable at every stage.

\section{Organisation of the Thesis}

The remainder of this thesis is organised as follows. Chapter~2 reviews the
related work in retrieval-augmented generation, agents, knowledge graphs and
document understanding, and identifies the gap that Synapse fills.
Chapter~3 presents the requirement analysis, including the use cases and the
functional and non-functional requirements. Chapter~4 describes the design and
architecture of the platform and the main data and control flows. Chapter~5
explains the implementation, covering the technology stack and the principal
modules. Chapter~6 presents the product itself through the web application, using
screenshots of the working system. Chapter~7 reports the testing and evaluation,
including the document benchmark and the purpose-built datasets. Chapter~8
concludes the thesis and outlines directions for future work.

\section{Summary}

This chapter introduced Synapse, a platform that lets a user understand,
question, visualise and map a private collection of documents through one web
application. It explained the motivation, namely that valuable knowledge is
locked inside mixed document collections and that general chatbots can neither
read such a collection nor answer about it honestly. It stated the problem, the
aim and the objectives, set the scope, listed the contributions, and described
the iterative methodology used to build the system. The next chapter reviews the
related work and shows where this project sits within it.
